% Options for packages loaded elsewhere
\PassOptionsToPackage{unicode}{hyperref}
\PassOptionsToPackage{hyphens}{url}
\documentclass[
]{article}
\usepackage{xcolor}
\usepackage{amsmath,amssymb}
\setcounter{secnumdepth}{-\maxdimen} % remove section numbering
\usepackage{iftex}
\ifPDFTeX
  \usepackage[T1]{fontenc}
  \usepackage[utf8]{inputenc}
  \usepackage{textcomp} % provide euro and other symbols
\else % if luatex or xetex
  \usepackage{unicode-math} % this also loads fontspec
  \defaultfontfeatures{Scale=MatchLowercase}
  \defaultfontfeatures[\rmfamily]{Ligatures=TeX,Scale=1}
\fi
\usepackage{lmodern}
\ifPDFTeX\else
  % xetex/luatex font selection
\fi
% Use upquote if available, for straight quotes in verbatim environments
\IfFileExists{upquote.sty}{\usepackage{upquote}}{}
\IfFileExists{microtype.sty}{% use microtype if available
  \usepackage[]{microtype}
  \UseMicrotypeSet[protrusion]{basicmath} % disable protrusion for tt fonts
}{}
\makeatletter
\@ifundefined{KOMAClassName}{% if non-KOMA class
  \IfFileExists{parskip.sty}{%
    \usepackage{parskip}
  }{% else
    \setlength{\parindent}{0pt}
    \setlength{\parskip}{6pt plus 2pt minus 1pt}}
}{% if KOMA class
  \KOMAoptions{parskip=half}}
\makeatother
\usepackage{longtable,booktabs,array}
\usepackage{multirow}
\usepackage{calc} % for calculating minipage widths
% Correct order of tables after \paragraph or \subparagraph
\usepackage{etoolbox}
\makeatletter
\patchcmd\longtable{\par}{\if@noskipsec\mbox{}\fi\par}{}{}
\makeatother
% Allow footnotes in longtable head/foot
\IfFileExists{footnotehyper.sty}{\usepackage{footnotehyper}}{\usepackage{footnote}}
\makesavenoteenv{longtable}
\usepackage{graphicx}
\makeatletter
\newsavebox\pandoc@box
\newcommand*\pandocbounded[1]{% scales image to fit in text height/width
  \sbox\pandoc@box{#1}%
  \Gscale@div\@tempa{\textheight}{\dimexpr\ht\pandoc@box+\dp\pandoc@box\relax}%
  \Gscale@div\@tempb{\linewidth}{\wd\pandoc@box}%
  \ifdim\@tempb\p@<\@tempa\p@\let\@tempa\@tempb\fi% select the smaller of both
  \ifdim\@tempa\p@<\p@\scalebox{\@tempa}{\usebox\pandoc@box}%
  \else\usebox{\pandoc@box}%
  \fi%
}
% Set default figure placement to htbp
\def\fps@figure{htbp}
\makeatother
\setlength{\emergencystretch}{3em} % prevent overfull lines
\providecommand{\tightlist}{%
  \setlength{\itemsep}{0pt}\setlength{\parskip}{0pt}}
\usepackage{bookmark}
\IfFileExists{xurl.sty}{\usepackage{xurl}}{} % add URL line breaks if available
\urlstyle{same}
\hypersetup{
  hidelinks,
  pdfcreator={LaTeX via pandoc}}

\author{}
\date{}

\begin{document}

\includegraphics[width=2.04167in,height=0.73819in]{media/rId20.png}

《软件工程》

\textbf{课程大作业}

\begin{longtable}[]{@{}
  >{\raggedleft\arraybackslash}p{(\linewidth - 2\tabcolsep) * \real{0.2954}}
  >{\centering\arraybackslash}p{(\linewidth - 2\tabcolsep) * \real{0.7046}}@{}}
\toprule\noalign{}
\endhead
\bottomrule\noalign{}
\endlastfoot
\textbf{组 号：} & \textbf{第三组} \\
\textbf{项目名称：} & \textbf{学校管理系统} \\
\textbf{专 业：} & \textbf{软件工程} \\
\textbf{班 级：} & \textbf{1603班} \\
\textbf{日 期：} & \textbf{2018年12月15日} \\
\end{longtable}

\textbf{要求：}

\begin{quote}
1、每组完成一个项目的软件工程文档一套，文档包含软件项目计划书，项目需求分析报告，项目的概要设计报告，项目的数据库设计报告、项目的详细设计报告、测试报告及项目总结。

2、每一个报告的内容按照文档模板，内容完整，内容详细，图表清晰且正确。

3、每一个报告的结构严格按照文档的模板结构搭建。

4、每一份报告的排版整齐，正确。一级标题三号宋体加粗，二级标题四号宋体加粗，三级标题小四号，宋体加粗，正文小四号，宋体。每一级标题设置段前段后0.5行，正文的行距固定值20磅。
\end{quote}

\textbf{赋分标准：}

1、以小组形式自拟题目，文档赋分标准如下：

\begin{longtable}[]{@{}
  >{\centering\arraybackslash}p{(\linewidth - 18\tabcolsep) * \real{0.0776}}
  >{\centering\arraybackslash}p{(\linewidth - 18\tabcolsep) * \real{0.0845}}
  >{\centering\arraybackslash}p{(\linewidth - 18\tabcolsep) * \real{0.0826}}
  >{\centering\arraybackslash}p{(\linewidth - 18\tabcolsep) * \real{0.1044}}
  >{\centering\arraybackslash}p{(\linewidth - 18\tabcolsep) * \real{0.0999}}
  >{\centering\arraybackslash}p{(\linewidth - 18\tabcolsep) * \real{0.1100}}
  >{\centering\arraybackslash}p{(\linewidth - 18\tabcolsep) * \real{0.0800}}
  >{\centering\arraybackslash}p{(\linewidth - 18\tabcolsep) * \real{0.1207}}
  >{\centering\arraybackslash}p{(\linewidth - 18\tabcolsep) * \real{0.1200}}
  >{\centering\arraybackslash}p{(\linewidth - 18\tabcolsep) * \real{0.1203}}@{}}
\toprule\noalign{}
\begin{minipage}[b]{\linewidth}\centering
\textbf{序号}
\end{minipage} & \begin{minipage}[b]{\linewidth}\centering
\textbf{1}
\end{minipage} & \begin{minipage}[b]{\linewidth}\centering
\textbf{2}
\end{minipage} & \begin{minipage}[b]{\linewidth}\centering
\textbf{3}
\end{minipage} & \begin{minipage}[b]{\linewidth}\centering
\textbf{4}
\end{minipage} & \begin{minipage}[b]{\linewidth}\centering
\textbf{5}
\end{minipage} & \begin{minipage}[b]{\linewidth}\centering
\textbf{6}
\end{minipage} & \begin{minipage}[b]{\linewidth}\centering
\textbf{7}
\end{minipage} & \begin{minipage}[b]{\linewidth}\centering
\textbf{8}
\end{minipage} & \begin{minipage}[b]{\linewidth}\centering
\textbf{9}
\end{minipage} \\
\midrule\noalign{}
\endhead
\bottomrule\noalign{}
\endlastfoot
\textbf{评分项} & 文档内容的完整性

（5\%） & 软件项目计划书

（15\%） & 项目需求分析报告

（20\%） & 项目概要

设计报告

（15\%） & 项目数据库

设计报告

（10\%） & 项目详细设计报告（15\%） & 项目测试报告

（10\%） & 项目总结报告

（5\%） & 项目团队的合作

（5\%） \\
\textbf{90-100} & 按照文档撰写格式，每项内容撰写充实、完整 &
项目计划书中内容完整，能够清晰看出项目的进度安排和成本核算，小组分工明确
&
报告内容完整，有清晰的项目功能的分析说明部分，并有结构图，有详细的数据输入输出说明及各种图形的表示
&
报告内容非常完整，报告中每个要素的内容充实且完整，各类图画的准确，流程清晰
&
报告内容非常完整，对数据分析清晰，数据概念设计中E-R图正确且关系清晰，各个表结构设计合理正确
&
报告内容完整，每个程序的描述及算法流程清晰准确，输入项及输出项内容描述正确。
& 报告内容完整，测试过程清晰，结论正确 &
站在不同角度说明本次项目总结的收获与存在问题，总结内容真实有深度，对存在问题说明真实
& 团队分工明确，证据真实且充实 \\
\textbf{80-90} & 按照文档撰写格式，每项内容都有撰写但不充实 &
项目计划书中内容较完整，较清晰看出项目的进度安排和成本核算，小组分工较明确
&
报告内容完整，项目功能的分析说明部分较清晰，并有结构图，有较详细的数据输入输出说明及各种图形的表示
&
报告内容非常完整，报告中每个要素的内容较充实完整，各类图画的较准确，流程较清晰
&
报告内容非常完整，对数据分析较清晰，数据概念设计中E-R图正确且关系较清晰，各个表结构设计较合理正确
&
报告内容完整，每个程序的描述及算法流程较清晰准确，输入项及输出项内容描述较正确。
& 报告内容完整，测试过程较清晰，结论正确 &
站在不同角度说明本次项目总结的收获与存在问题，总结内容真实较有深度，对存在问题总结较真实
& 团队分工明确，证据较真实，但不充实 \\
\textbf{70-80} &
按照文档撰写格式，只撰写部分内容，且撰写的内容比较充实， &
项目计划书中基本内容完整，较清晰看出项目的进度安排和成本核算，小组基本分工明确
& 报告内容基本完整，并有结构图，有数据输入输出说明基本基本及配图说明。 &
报告内容完整，报告中每个要素的内容基本充实完整，各类图画的基本准确，流程较清晰
&
报告内容较完整，对数据分析基本清晰，数据概念设计中E-R图正确且关系基本清晰，各个表结构设计基本合理正确
&
报告内容完整，每个程序的描述及算法流程基本清晰准确，输入项及输出项内容描述基本正确。
& 报告内容较完整，测试过程较清晰，结论基本正确 &
站在不同角度，基本说明本次项目总结的收获与存在问题，总结内容基本真实对存在问题说明基本真实
& 团队分工较明确，证据基本真实，但不充实 \\
\textbf{60-70} & 按照文档撰写格式，只撰写部分内容，且撰写的内容不充实 &
项目计划书中基本内容完整，有部分项目的进度安排和成本核算，无小组人员分工
& 有需求分析报告，有结构图 &
报告内容完整，报告中每个要素的内容完整，部分图画的准确，流程基本清晰 &
报告内容基本完整，数据概念设计中E-R图正确但关系不是很清晰，各个表结构设计存在欠缺
&
报告较内容完整，每个程序的描述及算法流程基本准确，输入项及输出项内容描述简单。
& 报告内容基本完整，测试基本过程清晰，结论基本正确 &
站在不同角度，基本说明本次项目总结的收获与存在问题 &
团队分工基本明确，证据基本真实 \\
\textbf{\textless60} & 没有文档按照格式撰写，内容也不充实 &
内容不完整，没有进度计划表 & 报告内容只有功能结构图，没有其他部分的内容
& 报告内容不完整 & 报告内容不完整，内容有错 &
报告内容不完整，且描述有错。 &
报告内容不完整，测试过程不清晰，结论不正确 & 总结内容应付 &
没有分工，没有证据 \\
\textbf{分值} & & & & & & & & & \\
\end{longtable}

2、以小组为单位提交文档资料后，个人成绩赋分标准如下表。

\begin{longtable}[]{@{}
  >{\centering\arraybackslash}p{(\linewidth - 10\tabcolsep) * \real{0.1935}}
  >{\centering\arraybackslash}p{(\linewidth - 10\tabcolsep) * \real{0.1361}}
  >{\centering\arraybackslash}p{(\linewidth - 10\tabcolsep) * \real{0.1429}}
  >{\centering\arraybackslash}p{(\linewidth - 10\tabcolsep) * \real{0.1746}}
  >{\centering\arraybackslash}p{(\linewidth - 10\tabcolsep) * \real{0.1865}}
  >{\centering\arraybackslash}p{(\linewidth - 10\tabcolsep) * \real{0.1665}}@{}}
\toprule\noalign{}
\begin{minipage}[b]{\linewidth}\centering
\textbf{评阅人}
\end{minipage} &
\multicolumn{5}{>{\centering\arraybackslash}p{(\linewidth - 10\tabcolsep) * \real{0.8065} + 8\tabcolsep}@{}}{%
\begin{minipage}[b]{\linewidth}\centering
\end{minipage}} \\
\midrule\noalign{}
\endhead
\bottomrule\noalign{}
\endlastfoot
\textbf{姓名} & \textbf{小组得分} & \textbf{小组人数} &
\textbf{小组成员总分} & \textbf{个人贡献度（\%）} & \textbf{个人成绩} \\
陈家乐 & \multirow{4}{=}{} & \multirow{4}{=}{4} & \multirow{4}{=}{} &
27\% & \\
张圆圆 & & & & 26\% & \\
谢犇 & & & & 24\% & \\
张钰 & & & & 23\% & \\
\end{longtable}

说明：以上总分为个人的期末成绩，本课程个人总成绩=终结性成绩*50\%+过程性成绩*50\%。

\textbf{工作量统计}

\begin{longtable}[]{@{}
  >{\centering\arraybackslash}p{(\linewidth - 18\tabcolsep) * \real{0.1471}}
  >{\centering\arraybackslash}p{(\linewidth - 18\tabcolsep) * \real{0.0781}}
  >{\centering\arraybackslash}p{(\linewidth - 18\tabcolsep) * \real{0.0955}}
  >{\centering\arraybackslash}p{(\linewidth - 18\tabcolsep) * \real{0.0954}}
  >{\centering\arraybackslash}p{(\linewidth - 18\tabcolsep) * \real{0.0954}}
  >{\centering\arraybackslash}p{(\linewidth - 18\tabcolsep) * \real{0.0954}}
  >{\centering\arraybackslash}p{(\linewidth - 18\tabcolsep) * \real{0.0954}}
  >{\centering\arraybackslash}p{(\linewidth - 18\tabcolsep) * \real{0.0954}}
  >{\centering\arraybackslash}p{(\linewidth - 18\tabcolsep) * \real{0.0954}}
  >{\centering\arraybackslash}p{(\linewidth - 18\tabcolsep) * \real{0.1072}}@{}}
\toprule\noalign{}
\begin{minipage}[b]{\linewidth}\centering
\textbf{姓名}
\end{minipage} & \begin{minipage}[b]{\linewidth}\centering
\textbf{软件}

\textbf{项目}

\textbf{计划书}
\end{minipage} & \begin{minipage}[b]{\linewidth}\centering
\textbf{需求}

\textbf{分析}

\textbf{报告}
\end{minipage} & \begin{minipage}[b]{\linewidth}\centering
\textbf{概要设计报告}
\end{minipage} & \begin{minipage}[b]{\linewidth}\centering
\textbf{数据库设计}

\textbf{报告}
\end{minipage} & \begin{minipage}[b]{\linewidth}\centering
\textbf{详细设计报告}
\end{minipage} & \begin{minipage}[b]{\linewidth}\centering
\textbf{系统测试报告}
\end{minipage} & \begin{minipage}[b]{\linewidth}\centering
\textbf{总结}

\textbf{报告}
\end{minipage} & \begin{minipage}[b]{\linewidth}\centering
\textbf{工作量}
\end{minipage} & \begin{minipage}[b]{\linewidth}\centering
\textbf{百分比（\%）}
\end{minipage} \\
\midrule\noalign{}
\endhead
\bottomrule\noalign{}
\endlastfoot
陈家乐 & 2 & 3 & 5 & 6 & 3 & 4 & 4 & 27 & 27\% \\
张圆圆 & 4 & 5 & 3 & 4 & 3 & 3 & 4 & 26 & 26\% \\
谢犇 & 3 & 4 & 3 & 3 & 4 & 4 & 3 & 24 & 24\% \\
张钰 & 3 & 2 & 4 & 3 & 5 & 3 & 3 & 23 & 23\% \\
\multicolumn{8}{@{}>{\centering\arraybackslash}p{(\linewidth - 18\tabcolsep) * \real{0.7974} + 14\tabcolsep}}{%
小组工作量合计} & 100 & 100\% \\
\end{longtable}

\end{document}
